% !TeX root = ../../Book3.tex
% 纲要标题：### B.4 模与自由消解的计算
% 纲要位置：工作记录/计划纲要.md:2853

\section{模与自由消解的计算}
\label{b3:sec:B04}

关系矩阵随生成元的更换而改变，计算记录必须保存这些更换的双向表示。尤其在消解中，相邻矩阵乘积为零只给出像包含于核，核的生成证书才证明正合性。

% 纲要内容（仅作写作计划，尚非教材正文）：
% * 调用第21章的 Schreyer 构造与分次极小化，以证书对照表区分全部核、复形、正合性与 Betti 数
% * 局部标准基方法归附录 F；本节证明局部成员与原环中乘性集倍数证书的等价性
% * SINGULAR 等精确代数系统可作实验载体；工作环、项序与局部化信息必须记录
% * 多项式与局部环中的同一输入可能有不同的理想成员答案，应提供对照实例
% #### B.4.1 Schreyer 构造、自由消解与 Betti 表
% #### B.4.2 局部计算中的项序与结果验证
% #### B.4.3 计算机代数系统中的环与局部化
% #### B.4.4 多项式环与局部环中理想成员判定的比较


\subsection{Schreyer 构造、自由消解与 Betti 表}
\label{b3:subsec:B04:01}

设 \(k\) 为域，\(S=k[x_1,\ldots,x_n]\)，输入矩阵 \(F\in S^{r\times s}\) 的列生成子模 \(N\subset S^r\)。计算模 Gröbner 基矩阵 \(G\in S^{r\times t}\) 时，保存双向表示
\[
G=FU,\qquad F=GV,
\quad U\in S^{s\times t},\quad V\in S^{t\times s}.
\]
由\cref{b3:thm:schreyer-syzygies}求出生成 \(\ker G\) 的矩阵 \(T\) 后，\cref{b3:alg:polynomial-matrix-kernel}将原核还原为
\[
\ker F=\operatorname{im}K,\qquad
K=\bigl[\,UT\ \ I_s-UV\,\bigr].
\]
第二块记录原生成组与新生成组之间多余的表示；它可能含有第一块没有给出的关系。

\begin{example}\label{b3:exm:B-generator-change}
设 \(k\) 为域，在 \(S=k[x,y]\) 中取 \(F=(x,y,x+y)\)，改用 \(G=(x,y)\)。可取
\[
U=\begin{pmatrix}1&0\\0&1\\0&0\end{pmatrix},
\quad
V=\begin{pmatrix}1&0&1\\0&1&1\end{pmatrix},
\quad T=\binom{y}{-x}.
\]
因此原关系由
\[
(y,-x,0)^{\mathsf T},\qquad(-1,-1,1)^{\mathsf T}
\]
生成。第二个关系来自被删掉的冗余生成元，单算新基的合冲而不记录变换，就会漏掉它。
\end{example}

同样的记录要逐层保存。\cref{b3:tab:B-resolution-certificates}列出各项结论所需的数据，其中 \(C_\bullet\xrightarrow{\varepsilon}M\) 是有限秩自由模的增广序列，\(d_i:C_i\rightarrow C_{i-1}\)；讨论分次极小性与 Betti 数时，\(S\) 取标准分次，\(M\) 有限生成，各映射齐次且次数为零。

\begin{table}[htbp]
\centering
\caption{合冲、消解与 Betti 数的计算证书}
\label{b3:tab:B-resolution-certificates}
\begin{tabularx}{\linewidth}{@{}>{\raggedright\arraybackslash}p{0.20\linewidth}>{\raggedright\arraybackslash}X@{}}
\toprule
所需结论 & 必须保存或核验的数据 \\
\midrule
生成子模相等 & \(G=FU\) 与 \(F=GV\) 的精确矩阵等式 \\
全部核已求出 & 已验证的模 Gröbner 基 \(G\)、Schreyer 合冲矩阵 \(T\) 及两向变换 \(U,V\)；由还原公式得到全部核生成矩阵 \(K\)，不只检查 \(FK=0\) \\
序列为复形 & 各层 \(d_i d_{i+1}=0\)，且 \(\varepsilon d_1=0\) \\
复形为消解 & \(\varepsilon\) 满射、\(\ker\varepsilon=\operatorname{im}d_1\)，以及每一层的全部核与下一层像相等；有限左端还须核验最左侧微分单射 \\
分次消解极小 & 已验证正合的消解，其微分项全在 \(\mathfrak m=(x_1,\ldots,x_n)\) 中；消去单位项时须同时更新相邻矩阵 \\
Betti 数 & 记录极小消解中 \(S(-j)\) 在 \(C_i\) 中的重数 \(\beta_{i,j}\)，它等于 \(\dim_k\operatorname{Tor}^S_i(M,k)_j\) \\
\bottomrule
\end{tabularx}
\end{table}

\cref{b3:thm:hilbert-syzygy}给出有限生成分次模具有有限长度自由消解的依据，\cref{b3:prop:graded-cancellation}给出单位项消去的具体步骤；极小消解与 Tor 的关系见\eqref{b3:eq:graded-betti-tor}。这些结论使求核与极小化各有明确的核验目标。

\subsection{局部计算中的项序与结果验证}
\label{b3:subsec:B04:02}

全局模项序的首项约化以良序为终止依据。局部序允许 \(1>x>x^2>\cdots\)，会改变可逆元素，也会改变约化方法。因此从多项式环改为局部计算时，不能仅把排序参数换掉，再沿用原来的除法终止证明。

\ProseParagraphBreak
\appref{b3:app:F}将定义局部标准基和弱正规形。局部成员关系可以完全写成原环中的有限等式。

\begin{proposition}[局部成员证书]\label{b3:prop:B-local-membership-certificate}
设 \(S\) 为交换环，\(U\subset S\) 为不含零的乘性集，\(R=U^{-1}S\)。给定 \(F\in S^{r\times s}\) 与 \(f\in S^r\)，则
\[
\dfrac{f}{1}\in\operatorname{im}(U^{-1}F:R^s\longrightarrow R^r)
\quad\Longleftrightarrow\quad
\text{存在 }u\in U,\ a\in S^s\text{，使 }uf=Fa.
\]
\end{proposition}
\begin{proof}
若 \(uf=Fa\)，在 \(R\) 中除以 \(u\) 即得成员表示。反过来，若 \(f/1\) 属于像，将表示中有限多个分母通分，可写成 \(f/1=Fb/v\)，其中 \(v\in U\)、\(b\in S^s\)。局部化中的相等意味着存在 \(w\in U\) 使 \(w(vf-Fb)=0\)。取 \(u=wv\)、\(a=wb\)，便有 \(uf=Fa\)。此处乘上 \(w\) 的步骤也涵盖 \(S\) 有零因子的情形。
\end{proof}

因此验证成员证书时，先检查原环中的矩阵等式，再核验 \(u\) 属于指定乘性集。在 \(S=k[x_1,\ldots,x_n]\)、\(U=S\setminus(x_1,\ldots,x_n)\) 的原点局部化中，后一个条件就是 \(u(0)\ne0\)。

\ProseParagraphBreak
非成员证明需要一个已验证的局部标准基及适用的正规形定理。仅有一次未约化为零的尝试不能作数，因为约化方法可能不完整。若输出为局部自由消解，每层仍要说明全部核的生成，并在作单位消去后再次检查相邻微分。局部极小性由极大理想中的微分项判断，而不是由多项式总次数或文件中列数的变化判断。

\subsection{计算机代数系统中的环与局部化}
\label{b3:subsec:B04:03}

一份可复算的环描述至少包括系数域、变量、变量次序、项序、商关系以及被倒置的元素。若系数是代数扩张，还要写出其定义多项式；若对象带分次或滤过，还要记录各变量和基向量的次数。只保存一张系数表无法恢复这些信息。

\ProseParagraphBreak
例如以下三个工作环有同名变量，却表达不同的问题：
\[
\mathbb Q[x,y],\qquad
\mathbb Q[x,y]_{(x,y)},\qquad
\mathbb Q[\![x,y]\!].
\]
第一环中的多项式 \(1-x\) 不可逆，后两个环中则可逆。第二环的元素是满足分母条件的有理函数；第三环还允许一般形式级数，所以它们并非相同的环。使用有限截断计算第三环时，还必须附上精度及为何该精度足够的说明。

\ProseParagraphBreak
SINGULAR 等精确计算系统可以承载这些实验，但报告不应以某个命令名称代替数学定义。应把“输入理想、输出标准基、表示矩阵、余项或 Hilbert 数据”分别保存。软件内部使用何种临界对策略可以更换；公开的结果仍应由明确的环、项序和有限等式决定。\subsecref{b3:subsec:F04:01}将给出截断何时成为完整长度证书的具体条件。

\subsection{多项式环与局部环中理想成员判定的比较}
\label{b3:subsec:B04:04}

\begin{example}\label{b3:exm:B-local-membership}
设 \(k\) 为域，\(S=k[x]\)，\(I=(x-x^2)\)。在 \(S\) 中，\(x\notin I\)：若 \(x=(x-x^2)h\)，整环中约去 \(x\) 得 \(1=(1-x)h\)，与 \(1-x\) 非单位矛盾。但在原点局部环 \(S_{(x)}\) 中，
\[
x=(1-x)^{-1}(x-x^2),
\]
故 \(IS_{(x)}=(x)\)。在点 \(x=1\) 的局部环中，反而是 \(x\) 可逆，于是 \(IS_{(x-1)}=(1-x)\)。
\end{example}

从商环看，这个差别同样清楚。多项式商 \(k[x]/(x-x^2)\) 有两个点，并由中国剩余定理同构于 \(k\times k\)；在原点局部化只保留一个分量。局部成员关系的改变对应实际删去远离所选点的部分，不是排序技巧造成的近似。

\ProseParagraphBreak
因此，计算前首先应决定研究的是整个代数集、指定点附近的局部环，还是完备局部环。全局求解可以先计算零维商代数，再按局部分量分析；局部标准基也可以直接计算指定点的结构。两种方法能够在适当比较映射下核对，但它们首先必须确实描述同一个局部对象。
